\section{Raw task order and changing-rank stability}\label{sec:order}
The central theorem is formulated in measures because selection acts before normalization. This section gives a verification principle from primitive kernels and a singular extension of that principle. It does not assume a covariance or information matrix in the experiment.

\subsection*{A measure determined by executable affine-square tests}
For a report--audit law $P(dy,dz)$ on $Y\times H$ define the finite positive matrix measure
\begin{equation}\label{eq:moment-matrix}
 J_P(D)=\int_{D\times H}v(z)v(z)^\mathsf T P(dy,dz),
 \qquad v(z)=(1,z_1,\ldots,z_k)^\mathsf T.
\end{equation}
For finite positive matrix measures write $A\preccurlyeq B$ when $B(D)-A(D)$ is positive semidefinite for every Borel $D$. This definition is independent of the choice of dominating scalar measure.

\begin{proposition}[An exact order criterion for localized affine-square tests]\label{prop:matrix-order}
For two joint laws $P,P^0$ on the same full report space and $0<l\le1\le u$, the following assertions are equivalent:
\begin{equation}\label{eq:matrix-order}
 lJ_{P^0}\preccurlyeq J_P\preccurlyeq uJ_{P^0};
\end{equation}
for every Borel report event $D$ and every $a\in\R^{k+1}$,
\begin{equation}\label{eq:affine-tests}
 l\int_D|a^\mathsf T v(z)|^2dP^0
 \le\int_D|a^\mathsf T v(z)|^2dP
 \le u\int_D|a^\mathsf T v(z)|^2dP^0.
\end{equation}
It suffices to test rational vectors $a$. These equivalent conditions imply, for all Borel $q\in[0,1]$ and finite-image Borel $r:Y\to H$, both the report-mass and the squared-task inequalities with zero additive defect and these same factors. A violation of \eqref{eq:matrix-order} is witnessed by a single Borel report event and a rational affine-square test.
\end{proposition}
\begin{proof}
The integral in \eqref{eq:affine-tests} is $a^\mathsf T J_P(D)a$. The definition of positive semidefiniteness gives the equivalence. Continuity of the quadratic form gives sufficiency of rational vectors and a rational witness for a strict violation. Taking $a=(1,0,\ldots,0)$ gives mass order. For a fixed readout $r$, sum the $k$ quadratic tests with vectors $(-r_j,e_j)$, obtaining the order of $\|z-r\|^2$. Partition $Y$ into the finitely many cells where $r$ is constant. For a simple nonnegative $q$, apply order on intersections with these cells and add. Bounded monotone approximation extends to arbitrary $q\in[0,1]$. All measures are finite and all tested losses bounded.
\end{proof}

Thus the criterion is necessary and sufficient for the declared task-rich affine-square family. It is a sufficient condition for the narrower squared-prediction task of the central theorem. It is not claimed necessary for that narrower task, or for the existence of an arbitrary causal balancing strategy. The matrix is built from raw joint audit laws, not from an assumed parameterization. The upper construction may use its virtual report space after averaging its current command; the lower criterion must use the full report seen by a competitor, conditional on that competitor's chosen command.

\begin{theorem}[Order on measurable singular ranges]\label{thm:singular-order}
Let $J_0,J_1$ be finite positive matrix measures of size $k+1$. With
$\lambda=\tr J_0+\tr J_1$, take measurable positive semidefinite densities $A_0,A_1$ with respect to $\lambda$. Then
\begin{equation}\label{eq:singular-equivalence}
 lJ_0\preccurlyeq J_1\preccurlyeq uJ_0
\end{equation}
is equivalent, $\lambda$-almost everywhere, to equality of kernels
$\ker A_0=\ker A_1$ and to
\begin{equation}\label{eq:range-spectrum}
 lP_0\preccurlyeq A_0^{\dagger/2}A_1A_0^{\dagger/2}
            \preccurlyeq uP_0,
\end{equation}
where $P_0$ is the orthogonal projection onto $\operatorname{ran}A_0$ and $\dagger$ denotes the Moore--Penrose inverse. On the zero range both sides are zero. All these maps are measurable, even when rank changes. No positive lower bound on a nonzero eigenvalue of $A_0$ is required. In particular every zero affine-square task on an acquired report event remains exactly zero under the order.
\end{theorem}
\begin{proof}
A finite signed matrix measure is positive if and only if its density is positive semidefinite almost everywhere. To see the nontrivial direction, test the density against all rational vectors; a negative set of positive measure for one such quadratic form contradicts positivity of its integral on that set. The countable intersection of the resulting full-measure sets handles all rational vectors, hence all real vectors by continuity. Thus \eqref{eq:singular-equivalence} is equivalent to $lA_0\preccurlyeq A_1\preccurlyeq uA_0$ almost everywhere.

For two such matrices, a vector in $\ker A_0$ has zero quadratic form under $A_1$, and positivity implies that it is in $\ker A_1$. The lower inequality gives the converse. On their common range $A_0$ is positive definite; congruence by its inverse square root yields \eqref{eq:range-spectrum}. Conversely equality of kernels makes $A_1$ supported on that range. Congruence of \eqref{eq:range-spectrum} by $A_0^{1/2}$ gives the original two matrix inequalities, extended by zero on the common kernel.

Measurability follows from finite-dimensional spectral functional calculus: $P_0$ is the pointwise limit of $A_0(A_0+n^{-1}I)^{-1}$, and the scalar function equal to $x^{-1/2}$ on $x>0$ and zero at zero is Borel. Applying it spectrally gives a Borel matrix function on each rank stratum, hence on their finite measurable union. We use these matrices pointwise; their unbounded inverse norms are never taken as a uniform error constant. Finally a nonnegative task of zero integral under $J_0$ has integral at most $u$ times zero under $J_1$, proving preservation of zero tasks.
\end{proof}

There is an important distinction between this assertion and a rank formula for risk. Rank is allowed to change; the rate still comes from the acquired masses and coverings of Section~\ref{sec:geometry}. For example a nonzero matrix density with arbitrarily small positive eigenvalues can obey the order with fixed $l,u$, while an arbitrarily small absolute perturbation can create a new positive direction in its kernel and fail the order. This is precisely the difference relevant to a zero-loss stratum.

\begin{proposition}[Positive-kernel verification and relative command calibration]\label{prop:positive-order}
Suppose raw positive input states satisfy $ls_0\le s\le us_0$ in an ordered vector space and a fixed positive normalized report--audit instrument sends an input to a joint probability law. Its output laws and task matrices satisfy the same order. More generally let $\mathcal I_j$ be positive instruments and let a nominal probability vector $(\lambda_j^0)$ satisfy a desired joint-law identity after averaging. If
\[
 (1-\eta)\lambda_j^0\le\lambda_j\le(1+\eta)\lambda_j^0
 \quad\hbox{for every }j,\qquad 0\le\eta<1,
\]
then the actual averaged joint law lies between $1-\eta$ and $1+\eta$ times the nominal law, for every input. Successive relative input and instrument bounds multiply their lower factors and upper factors. They do not add an acquisition-amplified defect.
\end{proposition}
\begin{proof}
Positivity sends the two positive differences $s-ls_0$ and $us_0-s$ to positive measures. Integration of $v(z)v(z)^\mathsf T$ against a positive measure preserves positive matrix order. The command assertion follows by multiplying each nonnegative instrument measure by its nonnegative weight difference and summing. Applying the two inequalities successively multiplies the factors. These arguments do not divide by any report probability.
\end{proof}

\subsection*{Actual-mass-weighted predictive stability}
The following is the trace/base-norm contraction principle for positive normalized instruments; in the quantum case compare \cite[Corollary 3.40]{Watrous18}. We include its proof to retain the unnormalized and zero-mass conventions used here.
\begin{proposition}[Continuity through likelihood zeros and pure states]\label{prop:continuity}
A classical stochastic instrument on signed measures, or a finite-dimensional completely positive normalized quantum instrument, maps normalized inputs $s,t$ to unnormalized post-report states $A_y,B_y$. Use any common finite dominating measure for these two outputs and let $p_y=u(A_y)$, $q_y=u(B_y)$ be their masses. Then
\begin{align}
 \int\|A_y-B_y\|_1dy&\le\|s-t\|_1,\label{eq:stable1}\\
 \|p-q\|_{\TV}&\le\tfrac12\|s-t\|_1,\label{eq:stable2}\\
 \int\min(p_y,q_y)\left\|\frac{A_y}{p_y}-\frac{B_y}{q_y}\right\|_1dy
 &\le2\|s-t\|_1.\label{eq:stable3}
\end{align}
Terms of zero mass are zero. The same unnormalized and report-law inequalities hold for a complete stopped instrument under one common legal policy which terminates almost surely under both inputs.
\end{proposition}
\begin{proof}
Write $s-t=D_+-D_-$ with positive parts. Positivity, the triangle inequality and preservation of integrated mass give
$\int\|\mathcal I_y(D_+)-\mathcal I_y(D_-)\|_1dy\le u(D_+)+u(D_-)=\|s-t\|_1$.
Taking masses gives \eqref{eq:stable2}. If $p\le q$, then
\[
 p\|A/p-B/q\|_1\le\|A-B\|_1+(q-p).
\]
Interchanging the roles for $q\le p$, integrating, and using $|p-q|\le\|A-B\|_1$ proves \eqref{eq:stable3}. A finite causal continuation with a fixed policy is again a positive normalized instrument. Truncate a stopped continuation at a deterministic horizon, retaining a distinct surviving outcome. Apply the finite inequality, then let the horizon increase. The surviving masses vanish under both inputs by almost-sure termination, and the sum over terminal times converges monotonically. No expected-time assertion is deduced from these norm inequalities.
\end{proof}

The upper finite-state construction never stores a normalized posterior. It uses the acquired report and a finite decoder. Consequently stability of the raw instrument, rather than a globally contracting posterior recursion, suffices for the singular realizations below. More general closed-loop optimal-control claims would require their own strategy comparison and stability; they are not consequences of comparing the same fixed policy.
